\documentclass{article}
\usepackage[utf8]{inputenc}
\usepackage{amsmath}
\usepackage{amssymb}
\usepackage{parskip}
\usepackage{csquotes}
\usepackage{microtype}
\usepackage{enumitem}
\usepackage{titlesec}
\usepackage{hyperref}

\title{A Manhattan Project for Collective Corrigibility}
\author{Albert Jan van Hoek}
\date{\today}

% Custom formatting for sections
\titleformat{\section}{\normalfont\Large\bfseries}{\thesection}{1em}{}
\titleformat{\subsection}{\normalfont\large\bfseries}{\thesubsection}{1em}{}

\begin{document}

\maketitle

\begin{center}
    \textbf{A global research and development programme for keeping human civilization capable of learning and collaborating under radical uncertainty}
\end{center}

\section{The Problem}
Human civilization faces threats whose nature, timing, and consequences we cannot reliably predict.

Some are already visible: pandemics, climate change, ecological degradation, geopolitical conflict, and rapidly developing technologies. Others may emerge from interactions between systems that we do not yet understand.

The common feature is not that we know nothing. It is that \textbf{our models of the world are necessarily incomplete, while the world itself continues to change}.

This creates a fundamental problem for collective decision-making. A decision is never made directly from reality. It is made from a model of reality. If the model is incomplete, the decision may be wrong. More importantly, we may not know that it is wrong.

The central question therefore cannot simply be:
\begin{quote}
    \textbf{How do we make the correct decision?}
\end{quote}

Under radical uncertainty, a more fundamental question is:
\begin{quote}
    \textbf{How do we remain capable of discovering that our current understanding and actions are wrong?}
\end{quote}

But there is a second problem. Collective learning does not occur between abstract models. It occurs between \textbf{people, institutions, and other agents that must remain connected long enough to learn from one another}.

\begin{itemize}
    \item If disagreement destroys relationships, correction becomes dangerous.
    \item If mistakes destroy trust, people conceal them.
    \item If conflict cannot be repaired, networks fragment.
    \item If relationships become too fragile, information stops flowing precisely when it is most needed.
\end{itemize}

The problem is therefore twofold:
\begin{quote}
    \textbf{How do we build collectives that can continuously correct their models?}
\end{quote}
and:
\begin{quote}
    \textbf{How do we build the relationships that allow those collectives to remain connected while doing so?}
\end{quote}

This proposal treats both as scientific and engineering problems.

\section{The Central Hypothesis}
The central hypothesis of this programme is:
\begin{quote}
    \textbf{A civilization's ability to survive radical uncertainty depends not only on the quality of its knowledge, but on its capacity to collectively detect, communicate, and correct errors while maintaining the relationships required for continued collaboration.}
\end{quote}

We call the first capability \textbf{collective corrigibility}. We call the second \textbf{sustainable collaboration}. Together they form the target capability.

\begin{itemize}
    \item A collectively corrigible system is not a system that is always right. It is a system that remains capable of becoming less wrong.
    \item A sustainably collaborative system is not one in which conflict is absent. It is one in which disagreement, error, and conflict can occur \textbf{without destroying the relationships needed for future learning}.
\end{itemize}

This distinction is fundamental. The objective is therefore not to create universal agreement about the world. It is to create a system in which \textbf{different imperfect models can continuously expose and correct one another's limitations while the relationships between their holders remain sufficiently strong for that process to continue}.

\section{The Innovation}
Most approaches to existential and global risks are threat-specific. They ask:
\begin{itemize}
    \item How do we prevent pandemics?
    \item How do we mitigate climate change?
    \item How do we govern artificial intelligence?
    \item How do we prevent nuclear war?
    \item How do we protect ecosystems?
\end{itemize}

These questions remain essential. But this programme asks a different question:
\begin{quote}
    \textbf{What capability would help us respond to all of these threats, including threats we have not yet identified?}
\end{quote}

The proposed innovation is therefore not a new solution to a particular threat. It is a \textbf{general-purpose capability for collective learning and sustainable collaboration under radical uncertainty}.

The intended outcome is a civilization that becomes better at:
\begin{enumerate}
    \item detecting when its models are failing;
    \item making relevant information visible;
    \item incorporating perspectives that differ from its own;
    \item correcting decisions;
    \item learning from failures;
    \item repairing relationships after disagreement and error;
    \item improving the processes by which it learns and collaborates;
    \item preserving these capabilities as systems scale.
\end{enumerate}

\section{The Object of Research}
The object of research is not primarily the individual. It is the \textbf{network of relationships between models}.

Every individual, institution, and scientific discipline observes only part of reality. Consequently, no participant can be expected to possess the complete model. This means that disagreement is not necessarily a failure of coordination. It may be information about the incompleteness of the collective model.

But disagreement also creates relational stress. The same interaction can therefore produce two very different outcomes:
\begin{itemize}
    \item \textbf{disagreement → information → correction → stronger understanding → stronger relationship}
    \item \textbf{disagreement → threat → defensive response → conflict → damaged relationship → information loss → model divergence}
\end{itemize}

The programme must therefore study both sides of the process. The key research question becomes:
\begin{quote}
    \textbf{Under what conditions do differences between models produce learning rather than fragmentation?}
\end{quote}
And:
\begin{quote}
    \textbf{Under what conditions can relationships absorb disagreement without suppressing it?}
\end{quote}

\section{The First Principle: Epistemic Limitation}
The foundation is recognition of a simple fact:
\begin{quote}
    \textbf{Every participant has limited access to reality.}
\end{quote}

This includes individuals, organizations, scientific disciplines, governments, and technological systems. The first capability required is therefore not certainty. It is the ability to recognize the boundaries of one's own model.

Participants should be able to distinguish:
\begin{itemize}
    \item \textbf{Observation}: What did I actually observe?
    \item \textbf{Interpretation}: What did I infer from that observation?
    \item \textbf{Assumption}: What am I assuming to be true?
    \item \textbf{Prediction}: What do I expect to happen?
    \item \textbf{Uncertainty}: What might I be wrong about?
    \item \textbf{Blind spot}: What could another observer see that I cannot?
\end{itemize}

This creates the epistemic foundation for collaboration.

\section{The Second Principle: The Verbs of Learning}
The programme would identify and develop a set of observable behaviours that enable correction. These include:
\begin{itemize}
    \item observing;
    \item distinguishing observation from interpretation;
    \item questioning assumptions;
    \item listening;
    \item explaining;
    \item challenging;
    \item seeking disconfirming evidence;
    \item expressing uncertainty;
    \item correcting;
    \item accepting correction;
    \item updating;
    \item integrating new information;
    \item supporting another person's learning;
    \item preserving connection despite disagreement.
\end{itemize}

These behaviours should not be treated merely as moral virtues. They can be understood as \textbf{functional components of a distributed learning system}. For example:
\begin{itemize}
    \item \textbf{Honesty} protects information fidelity.
    \item \textbf{Listening} allows information to cross model boundaries.
    \item \textbf{Questioning} exposes assumptions.
    \item \textbf{Correction} generates error signals.
    \item \textbf{Curiosity} drives exploration.
    \item \textbf{Uncertainty} identifies where additional information may be valuable.
    \item \textbf{Willingness to update} allows models to change.
    \item \textbf{Forgiveness} preserves the relationships through which future information can travel.
\end{itemize}

\section{The Third Principle: Relationships Are Infrastructure}
This principle should be treated as fundamental. A learning network requires connections. Those connections are not merely communication channels. They are relationships between agents who must remain willing and able to exchange information, challenge one another, and cooperate again in the future.

Relationship strength should therefore be considered a \textbf{system property}. A strong relationship does not mean that two participants agree. It means that they retain sufficient capacity to:
\begin{itemize}
    \item disagree;
    \item communicate honestly;
    \item tolerate uncertainty;
    \item correct one another;
    \item experience conflict;
    \item repair conflict;
    \item forgive mistakes;
    \item reconnect;
    \item and continue collaborating.
\end{itemize}

This leads to an important distinction:
\begin{quote}
    \textbf{Trust is not the absence of disagreement.}
\end{quote}

A stronger form of trust may be:
\begin{quote}
    \textbf{Confidence that the relationship can survive disagreement.}
\end{quote}

That property is particularly important in a system designed for correction.

\section{Conflict Is Not Necessarily System Failure}
A system that attempts to eliminate conflict may accidentally eliminate information. If different models are genuinely different, disagreement is inevitable.

The objective should therefore not be:
\begin{quote}
    \textbf{Prevent conflict.}
\end{quote}
It should be:
\begin{quote}
    \textbf{Maintain the capacity to have productive conflict and restore collaboration afterwards.}
\end{quote}

Conflict becomes a process with at least three stages:
\begin{itemize}
    \item \textbf{Divergence}: Our models or interests differ.
    \item \textbf{Engagement}: We expose and examine the difference.
    \item \textbf{Repair}: We restore sufficient relational capacity to continue working together.
\end{itemize}

A mature collective therefore needs not only \textbf{conflict capacity}, but \textbf{conflict recovery capacity}.

\section{Forgiveness as a System Property}
Forgiveness deserves explicit attention because it has a functional role in a learning network. Human beings inevitably make mistakes. If every mistake permanently damages a relationship, participants will eventually become reluctant to:
\begin{itemize}
    \item experiment;
    \item admit uncertainty;
    \item disclose errors;
    \item challenge others;
    \item take intellectual risks.
\end{itemize}

The network becomes increasingly conservative and increasingly opaque. Forgiveness creates the possibility of \textbf{re-entry after error}. But forgiveness should not mean ignoring harm or abandoning accountability.

The relevant property is:
\begin{quote}
    \textbf{Can a relationship recover from an error sufficiently quickly to restore its capacity for collaboration and learning?}
\end{quote}

This suggests that forgiveness itself may be measurable. Potential measures include:
\begin{itemize}
    \item time from conflict to attempted repair;
    \item probability of successful repair;
    \item time until collaboration resumes;
    \item whether relevant information continues to flow during conflict;
    \item whether the conflict produces learning;
    \item whether the same conflict recurs;
    \item whether trust is restored;
    \item whether the relationship becomes stronger, weaker, or merely functional after repair.
\end{itemize}

We might therefore define a concept such as:
\begin{quote}
    \textbf{relational recovery time}
\end{quote}
the time required for a relationship to regain sufficient functional capacity following conflict or error. This is analogous to recovery time in other resilient systems.

\section{Relationship Strength as a Measurable Variable}
The programme should therefore not measure only knowledge and decision quality. It should measure the state of the relationships connecting the system.

Potential dimensions include:
\begin{itemize}
    \item \textbf{Connection}: Do participants remain connected?
    \item \textbf{Information permeability}: Can relevant information cross the relationship?
    \item \textbf{Correction safety}: Can participants challenge one another?
    \item \textbf{Conflict capacity}: Can disagreement occur without immediate relational collapse?
    \item \textbf{Repair capacity}: Can damaged relationships recover?
    \item \textbf{Recovery speed}: How quickly can collaboration resume?
    \item \textbf{Forgiveness capacity}: Can mistakes be incorporated without permanently excluding the participant?
    \item \textbf{Reciprocity}: Do participants both give and receive correction?
    \item \textbf{Persistence}: Can relationships survive repeated disagreement?
    \item \textbf{Trust resilience}: Does trust recover after violations?
\end{itemize}

These properties could become part of a broader \textbf{relationship-strength index}.

\section{The Learning Loop Becomes Relational}
The basic learning loop is:
\begin{quote}
    \textbf{Model → Action → Observation → Error → Update → New Model}
\end{quote}

But in a human collective, this is incomplete. We need:
\begin{quote}
    \textbf{Model → Interaction → Observation → Disagreement → Conflict → Repair → Learning → Update → New Model}
\end{quote}

And simultaneously:
\begin{quote}
    \textbf{How did we learn?}
\end{quote}
and:
\begin{quote}
    \textbf{How did the relationship survive the learning process?}
\end{quote}

This creates a second-order recursive structure. The system learns about the world. It learns how it learns about the world. And it learns how to maintain the relationships through which that learning occurs.

\section{The Dual Optimization Problem}
This creates an important design constraint. We must avoid optimizing learning at the expense of relationships. And we must avoid optimizing harmony at the expense of learning.

\begin{itemize}
    \item A system in which everyone gets along but nobody can challenge anyone is not collectively intelligent.
    \item A system in which everyone constantly attacks one another but nobody can maintain a relationship is not collectively intelligent either.
\end{itemize}

The target lies between these failures:
\begin{quote}
    \textbf{Maximum capacity for correction while preserving sufficient relational stability for continued collaboration.}
\end{quote}

This should itself become an empirical research question.

\section{The Research Programme}
The programme should proceed through a sequence of empirical stages.

\subsection{Phase I — Define}
Develop a formal theory of:
\begin{itemize}
    \item collective corrigibility;
    \item sustainable collaboration;
    \item relationship strength;
    \item conflict capacity;
    \item relational recovery;
    \item forgiveness;
    \item information flow;
    \item model diversity;
    \item collective learning.
\end{itemize}

A central objective should be to define the relationships between these variables. For example:
\begin{itemize}
    \item Does stronger relationship resilience increase the willingness to expose uncertainty?
    \item Does faster conflict recovery increase information flow?
    \item Does excessive relational harmony reduce useful disagreement?
    \item Is there an optimal balance between model independence and relational connectivity?
\end{itemize}

\subsection{Phase II — Measure}
Develop instruments capable of measuring both collective learning and relational resilience.

\subsubsection{Epistemic Performance}
\begin{itemize}
    \item error detection;
    \item correction speed;
    \item model updating;
    \item information flow;
    \item model diversity;
    \item prediction accuracy;
    \item learning from failure;
    \item meta-learning.
\end{itemize}

\subsubsection{Relational Performance}
\begin{itemize}
    \item relationship strength;
    \item willingness to communicate uncertainty;
    \item conflict capacity;
    \item repair probability;
    \item relational recovery time;
    \item forgiveness/re-entry;
    \item persistence of collaboration;
    \item information flow after conflict.
\end{itemize}

The key innovation is that these should not be treated as separate outcomes. We should study their interaction.

\subsection{Phase III — Experiment}
Construct controlled environments in which different collective architectures can be compared. Groups could receive identical problems but different communication and relationship structures.

Researchers could introduce:
\begin{itemize}
    \item incomplete information;
    \item conflicting evidence;
    \item unexpected events;
    \item misleading information;
    \item changing environments;
    \item asymmetric information;
    \item model disagreement;
    \item interpersonal conflict;
    \item errors by group members.
\end{itemize}

The key outcomes would include both:
\begin{quote}
    \textbf{How quickly does the group discover and correct its mistakes?}
\end{quote}
and:
\begin{quote}
    \textbf{How well does the group preserve its ability to collaborate afterwards?}
\end{quote}

This would allow us to test whether relational resilience improves epistemic performance.

\subsection{Phase IV — Intervention}
Once measurable mechanisms have been identified, develop interventions. These might include:
\begin{itemize}
    \item structured disagreement;
    \item uncertainty protocols;
    \item correction protocols;
    \item conflict-repair procedures;
    \item forgiveness and re-entry mechanisms;
    \item reflective practices;
    \item peer review;
    \item prediction and retrospective evaluation;
    \item dissent channels;
    \item relationship maintenance practices;
    \item institutional safeguards for error reporting;
    \item training in listening, correction, and updating.
\end{itemize}

Every intervention should be treated as a hypothesis. For every intervention we should ask:
\begin{quote}
    Does it increase collective corrigibility?
\end{quote}
\begin{quote}
    Does it strengthen sustainable collaboration?
\end{quote}
\begin{quote}
    Does it improve both without sacrificing either?
\end{quote}

If not, it should be discarded or redesigned.

\subsection{Phase V — Real-World Pilots}
The next step is deployment in real institutions. Potential environments include:
\begin{itemize}
    \item scientific research teams;
    \item hospitals;
    \item public-health organizations;
    \item schools;
    \item companies;
    \item governments;
    \item international collaborations;
    \item online communities.
\end{itemize}

Different environments provide different combinations of uncertainty, hierarchy, conflict, and interdependence. This allows the programme to determine whether the underlying principles generalize across contexts.

\subsection{Phase VI — Stress Testing}
A civilization-scale system must be tested under conditions resembling the situations in which it is most needed. This means deliberately introducing:
\begin{itemize}
    \item rapidly changing conditions;
    \item conflicting evidence;
    \item misinformation;
    \item strategic deception;
    \item polarization;
    \item resource scarcity;
    \item time pressure;
    \item repeated mistakes;
    \item interpersonal conflict;
    \item institutional incentives;
    \item cascading failures.
\end{itemize}

The key question is not simply:
\begin{quote}
    Can the system make good decisions under stress?
\end{quote}
It is:
\begin{quote}
    \textbf{Can the system preserve its capacity to learn and remain connected under stress?}
\end{quote}

A system that becomes incapable of correction precisely when relationships become strained is not resilient.

\subsection{Phase VII — Network Experiments}
The next level is multiple interconnected collectives. Can independent learning communities correct one another? Can they maintain relationships across institutional and cultural boundaries? Can they disagree without fragmenting? Can they repair relationships after serious disagreement? Can decentralized systems collectively learn without requiring one central authority to possess the correct model?

This is where the architecture becomes civilization-scale. The objective is not to create one universal model. It is to create a network of \textbf{partially independent, mutually corrigible, and relationally resilient models}.

\subsection{Phase VIII — Institutionalization}
Successful mechanisms would be incorporated into institutions. Scientific institutions already institutionalize peer review. Engineering institutions institutionalize safety testing. Public-health systems institutionalize surveillance. Democratic systems institutionalize mechanisms for changing leadership.

A mature civilization could institutionalize something broader:
\begin{quote}
    \textbf{continuous correction of the models underlying collective decisions, together with continuous maintenance of the relationships required to perform that correction.}
\end{quote}

\subsection{Phase IX — Civilization-Scale Deployment}
Only after the preceding stages demonstrate measurable benefits should the programme attempt civilization-scale deployment.

The output would not be a single centralized organization. It would be:
\begin{itemize}
    \item principles;
    \item protocols;
    \item measurement tools;
    \item training systems;
    \item conflict-repair mechanisms;
    \item institutional designs;
    \item open standards;
    \item and methods for evaluating collective learning.
\end{itemize}

Different communities could implement them locally. The commonality would not be the answer. It would be the \textbf{capacity to learn together without destroying the relationships through which learning occurs}.

\section{The Success Criterion}
The programme needs a strict definition of success. Success is not:
\begin{itemize}
    \item universal agreement;
    \item absence of conflict;
    \item permanent harmony;
    \item perfect prediction;
    \item elimination of mistakes;
    \item creation of one correct model.
\end{itemize}

Success is:
\begin{quote}
    \textbf{A measurable increase in the ability of a collective to detect, communicate, correct, and learn from errors while maintaining the relationships required for continued collaboration.}
\end{quote}

The ultimate system should therefore be evaluated on two interacting dimensions:
\begin{itemize}
    \item \textbf{Collective corrigibility}: Can we become less wrong?
    \item \textbf{Sustainable collaboration}: Can we remain connected long enough to keep learning?
\end{itemize}

The civilization we are trying to build must possess both.

\section{The Ultimate Resilience Metric}
This suggests a deeper possibility. The resilience of a collective may depend not only on how quickly it can correct an error, but on whether it can \textbf{recover its capacity for correction after relational damage}.

A useful conceptual quantity might therefore be:
\begin{quote}
    \textbf{Time to collective recovery}
\end{quote}
The time required for a network to return to effective information exchange, disagreement, and joint learning after a significant disturbance.

This disturbance could be:
\begin{itemize}
    \item an incorrect decision;
    \item a betrayal of trust;
    \item a major disagreement;
    \item institutional failure;
    \item public failure;
    \item political conflict;
    \item or an unexpected external shock.
\end{itemize}

A resilient collective is therefore not one that avoids disturbance. It is one that can \textbf{return to learning after disturbance}.

\section{The Deeper Principle}
This leads to a broader understanding of collaboration. Collaboration is not simply people working together toward the same goal. It is the maintenance of a relationship in which:
\begin{quote}
    \textbf{I can tell you that I think you are wrong.}
\end{quote}
and:
\begin{quote}
    \textbf{You can tell me that I am wrong.}
\end{quote}
and:
\begin{quote}
    \textbf{We can survive discovering that one of us was wrong.}
\end{quote}
and eventually:
\begin{quote}
    \textbf{We can use what we discovered to become better together.}
\end{quote}

That is a much stronger definition of sustainable collaboration.

\section{The Civilization-Scale Objective}
Human civilization has become extraordinarily powerful at changing the world. Our ability to act increasingly exceeds our ability to predict the consequences of our actions. That creates a growing asymmetry:
\begin{quote}
    \textbf{more power + incomplete models + interconnected systems = greater consequences of error.}
\end{quote}

The answer cannot be to wait until we possess a complete model. We never will. The alternative is to develop a civilization capable of:
\begin{itemize}
    \item \textbf{continuous model correction}
    \item \textbf{continuous relationship repair.}
\end{itemize}

The two are connected. We need relationships strong enough to permit correction. And we need correction mechanisms strong enough to prevent relationships from becoming dependent on agreement.

\section{The Manhattan Project Analogy}
The original Manhattan Project demonstrated that humanity can mobilize extraordinary intellectual, institutional, and material resources around a problem considered strategically essential.

The proposal here is to apply that capacity to a fundamentally different problem. Not:
\begin{quote}
    \textbf{How do we solve the next threat?}
\end{quote}
But:
\begin{quote}
    \textbf{How do we build the capacity to solve threats we cannot yet define?}
\end{quote}

And not merely:
\begin{quote}
    \textbf{How do we learn together?}
\end{quote}
But:
\begin{quote}
    \textbf{How do we remain capable of learning together, even when learning creates disagreement, conflict, and error?}
\end{quote}

The project would therefore have no single technological endpoint. Its endpoint would be a new civilizational capability:
\begin{quote}
    \textbf{A civilization that can remain collectively corrigible and relationally resilient as the world changes.}
\end{quote}

\section{The First Concrete Experiment}
We do not need to begin by convincing humanity. We need to test the hypothesis.

The first objective should therefore be modest:
\begin{quote}
    \textbf{Build small experimental collectives, define measurable criteria for collective corrigibility and sustainable collaboration, introduce controlled uncertainty, model disagreement and relational conflict, and determine whether specific practices improve both the speed and quality of collective correction and the ability of relationships to recover.}
\end{quote}

We should deliberately create situations in which:
\begin{itemize}
    \item someone is wrong;
    \item someone discovers they are wrong;
    \item people disagree;
    \item relationships become strained;
    \item repair becomes necessary;
    \item and the collective must learn afterwards.
\end{itemize}

The goal is not to prevent these events. The goal is to understand \textbf{what allows a collective to pass through them without losing its capacity to learn}.

If a proposed intervention fails, we learn. If it succeeds, we test whether the effect survives increasing complexity. Then we scale.

\section{The Principle}
The entire programme can ultimately be reduced to two propositions:
\begin{quote}
    \textbf{We cannot guarantee that humanity will always understand the world correctly.}
\end{quote}
\begin{quote}
    \textbf{We can attempt to build a civilization that remains capable of discovering when it does not — and remains capable of doing so together.}
\end{quote}

That is the proposed innovation. Not a solution to climate change. Not a solution to pandemics. Not a solution to AI. Not a solution to war. But a general-purpose capability that may improve our ability to confront \textbf{all of them — and the threats we have not yet imagined}.

The goal is not to make humanity permanently right. The goal is not even to make humanity permanently harmonious.

\textbf{The goal is to make humanity remain capable of becoming less wrong, without losing the relationships that make collective learning possible.}

Or, more simply:
\begin{quote}
    \textbf{We need a civilization that can disagree, repair, learn, and continue.}
\end{quote}

\end{document}